\chapter{Orbitais de fragmentos e moléculas poliatômicas}\label{ch:b2:fragment-orbitals}

A estrutura de Lewis da água mostra duas ligações O--H equivalentes e dois
pares não ligantes equivalentes. No entanto, quando a água é ionizada por luz
ultravioleta, os elétrons saem com quatro energias diferentes, e não duas; o
metano, com quatro ligações C--H equivalentes, mostra duas. Os elétrons de uma
molécula não ficam em ligações entre pares de átomos: eles ocupam orbitais
estendidos sobre a molécula inteira. Este capítulo constrói esses orbitais a
partir dos \omterm{def:b2:fragment-orbitals:fragment}{orbitais de fragmentos}, explica por que a água é angular e o eteno
é plano, e por que um carbocátion prefere ter grupos alquila à sua volta.

\begin{recall}
\Cref{ch:b2:diatomic-mos}: \omterm{def:b2:diatomic-mos:lcao}{CLOA}, sobreposição e simetria, dois níveis em
interação, \omterm{def:b2:diatomic-mos:bonding}{orbitais ligantes} e antiligantes. O volume do 1.º ano: as formas
VSEPR, a hibridização como descrição, a ressonância, os carbocátions e sua
ordem de estabilidade (a interação das ligações C--H com o orbital vazio foi
anunciada ali).
\end{recall}

\section{O método dos fragmentos}

\begin{definition}[Fragmentos]\label{def:b2:fragment-orbitals:fragment}
Um \emph{fragmento}\index{fragmento} é um grupo de átomos de uma molécula
considerado isoladamente (um átomo, um par de hidrogênios, um grupo \ce{CH2}), e
um \emph{orbital de fragmento}\index{orbital de fragmento} é um de seus orbitais.
Os orbitais da molécula são construídos fazendo interagir os orbitais de dois
fragmentos, dois a dois, como os orbitais atômicos de uma molécula diatômica.
\end{definition}

\begin{definition}[Combinação adaptada à simetria]\label{def:b2:fragment-orbitals:symmetry-adapted}
Uma \emph{combinação adaptada à simetria}\index{combinacao adaptada a simetria@combinação adaptada à simetria} de
orbitais atômicos equivalentes é uma combinação que cada operação de simetria
da molécula (uma reflexão em um plano de simetria, uma rotação em torno de um
eixo) deixa inalterada ou transforma em sua oposta, ou, para conjuntos
degenerados, transforma em combinações do mesmo conjunto.
\end{definition}

\begin{proposition}[Dois hidrogênios]\label{prop:b2:fragment-orbitals:h2-pair}
Para dois hidrogênios equivalentes com orbitais 1s $h_1$ e $h_2$, as
\omterm{def:b2:fragment-orbitals:symmetry-adapted}{combinações adaptadas à simetria} são $h_1 + h_2$ e $h_1 - h_2$: a primeira não
muda, a segunda muda de sinal, sob qualquer operação que troque os dois
átomos.
\end{proposition}

\begin{proof}
Uma operação que troca os átomos transforma $h_1$ em $h_2$ e $h_2$ em
$h_1$: $h_1 + h_2$ não muda e $h_1 - h_2$ se torna $h_2 - h_1 = -(h_1 - h_2)$.
\end{proof}

\begin{proposition}[Três e quatro hidrogênios]\label{prop:b2:fragment-orbitals:h4}
Para quatro hidrogênios nos vértices de um tetraedro, as combinações são uma
combinação totalmente simétrica $h_1 + h_2 + h_3 + h_4$ e um conjunto de três
combinações degeneradas com um plano nodal cada uma, com a forma de orbitais
p. Para três hidrogênios em triângulo (como na amônia), uma combinação
simétrica e um par degenerado.
\end{proposition}

\admitted

As formas das combinações degeneradas, e seus nomes (a, e, t), vêm da teoria
de grupos, desenvolvida no volume do 3.º ano; aqui os nomes são apenas rótulos.

\begin{proposition}[Só o semelhante interage com o semelhante]\label{prop:b2:fragment-orbitals:interaction-rules}
Um \omterm{def:b2:fragment-orbitals:fragment}{orbital de fragmento} só interage com os orbitais do outro \omterm{def:b2:fragment-orbitals:fragment}{fragmento} que se
comportam da mesma maneira sob todas as operações de simetria da molécula; a
força da interação cresce com a sobreposição e diminui com a diferença de
energia.
\end{proposition}

\begin{proof}
Se dois orbitais se comportam de modo diferente sob alguma reflexão, um é
simétrico e o outro antissimétrico em relação a esse plano, e sua sobreposição
e sua \omterm{def:b2:diatomic-mos:integrals}{integral de ressonância} se anulam (\cref{prop:b2:diatomic-mos:symmetry}).
A dependência em relação à sobreposição e à diferença de energia é a do
\cref{thm:b2:diatomic-mos:heteronuclear}.
\end{proof}

\begin{method}[Um diagrama de fragmentos]\label{met:b2:fragment-orbitals:fragment-diagram}
\begin{enumerate}
  \item Divida a molécula em dois \omterm{def:b2:fragment-orbitals:fragment}{fragmentos}, em geral um átomo central e o
        conjunto dos átomos à sua volta.
  \item Construa as \omterm{def:b2:fragment-orbitals:symmetry-adapted}{combinações adaptadas à simetria} dos orbitais externos.
  \item Classifique os orbitais dos dois \omterm{def:b2:fragment-orbitals:fragment}{fragmentos} por seu comportamento sob
        os planos de simetria e os eixos da molécula.
  \item Faça cada orbital interagir com os de mesmo comportamento, mais
        fortemente quando próximos em energia; os que sobram permanecem não
        ligantes.
  \item Preencha com os elétrons de valência e leia: nível ocupado mais alto,
        caráter ligante, preferências de forma.
\end{enumerate}
\end{method}

\section{A água}

Coloque a água no plano $yz$, com o eixo $z$ na bissetriz do ângulo H--O--H.
Dois planos de simetria contêm o eixo $z$: o plano molecular $yz$ e o plano
$xz$, perpendicular a ele. Em relação ao plano $xz$, que troca os dois
hidrogênios, $h_1 + h_2$ é simétrica e $h_1 - h_2$ antissimétrica. No oxigênio, os
orbitais 2s e $2\mathrm p_z$ são simétricos em relação aos dois planos; o
$2\mathrm p_y$ (no plano molecular, perpendicular a $z$) é antissimétrico em
relação a $xz$; o $2\mathrm p_x$ é antissimétrico em relação ao plano
molecular, o que nenhuma combinação dos hidrogênios é.

Assim, $h_1 + h_2$ se mistura com 2s e $2\mathrm p_z$, $h_1 - h_2$ com $2\mathrm p_y$, e
$2\mathrm p_x$ permanece não ligante. Resultam seis orbitais, quatro deles
preenchidos pelos oito elétrons de valência.

\begin{omfigure}
\begin{tikzpicture}[font=\small]
  % O fragment (left)
  \pic (os) at (0,-2.4) {omlevel={ud}{0.7}}; \node[left] at (-0.45,-2.4) {2s};
  \pic (op1) at (-0.45,0) {omlevel={ud}{0.4}}; \pic (op2) at (0,0) {omlevel={u}{0.4}}; \pic (op3) at (0.45,0) {omlevel={u}{0.4}};
  \node[left] at (-0.7,0) {2p};
  \node at (0,-3.2) {O};
  % H2 fragment (right)
  \pic (hp) at (6,0.9) {omlevel={u}{0.7}}; \node[right] at (6.4,0.9) {$h_1 + h_2$};
  \pic (hm) at (6,1.8) {omlevel={u}{0.7}}; \node[right] at (6.4,1.8) {$h_1 - h_2$};
  \node[anchor=north west] at (6.4,0.6) {fragmento \ce{H2}};
  % MOs
  \pic (m1) at (3,-2.9) {omlevel={ud}{0.7}};  \node[omsymlabel, below=0.17cm, inner sep=0.5pt] at (m1-c) {$2a_1$};
  \pic (m2) at (3,-1.55) {omlevel={ud}{0.7}};  \node[omsymlabel, below=0.17cm, inner sep=0.5pt] at (m2-c) {$1b_2$};
  \pic (m3) at (3,-0.78) {omlevel={ud}{0.7}};  \node[omsymlabel, below=0.17cm, inner sep=0.5pt] at (m3-c) {$3a_1$};
  \pic (m4) at (3,0) {omlevel={ud}{0.7}};  \node[omsymlabel, below=0.17cm, inner sep=0.5pt] at (m4-c) {$1b_1$};
  \pic (m5) at (3,2.4) {omlevel={empty}{0.7}}; \node[omsymlabel, below=0.17cm, inner sep=0.5pt] at (m5-c) {$4a_1$};
  \pic (m6) at (3,3.1) {omlevel={empty}{0.7}}; \node[omsymlabel, below=0.17cm, inner sep=0.5pt] at (m6-c) {$2b_2$};
  \draw[omcorr, densely dashed] (os-r) -- (m1-l) (op3-r) -- (m2-l) (op3-r) -- (m3-l) (op3-r) -- (m4-l) (op3-r) -- (m5-l) (op3-r) -- (m6-l);
  \draw[omcorr, densely dashed] (hp-l) -- (m1-r) (hm-l) -- (m2-r) (hp-l) -- (m3-r) (hp-l) -- (m5-r) (hm-l) -- (m6-r);
\end{tikzpicture}

{\small Diagrama de \omterm{def:b2:fragment-orbitals:fragment}{fragmentos} da água, esquemático. A combinação simétrica
dos hidrogênios se mistura com os orbitais 2s e $2\mathrm p_z$ do oxigênio
(orbitais $a_1$), a antissimétrica com $2\mathrm p_y$ ($b_2$); o $2\mathrm p_x$,
perpendicular ao plano molecular, permanece não ligante ($1b_1$, no nível do
2p do oxigênio). Quatro níveis de valência preenchidos, quatro bandas de
fotoelétrons; o mais alto, $1b_1$, é ionizado a \qty{12.62}{eV}. Os rótulos
$a_1$, $b_1$, $b_2$ são nomes (vindos da teoria de grupos, no volume do 3.º ano).}
% ledger: ie:H2O
\end{omfigure}

\begin{proposition}[Por que a água é angular]\label{prop:b2:fragment-orbitals:bent-water}
Em um H--O--H linear, dois orbitais p do oxigênio perpendiculares ao eixo são
não ligantes. A flexão permite a um deles, o que está no plano molecular, se
sobrepor a $h_1 + h_2$ e se misturar com o orbital 2s: esse nível preenchido
desce, mais do que sobem os outros níveis preenchidos. Com oito elétrons de
valência, a molécula angular é a mais estável.
\end{proposition}

\begin{proof}[Raciocínio]
Na molécula linear ($z$ ao longo de H--O--H), $h_1 + h_2$ tem sobreposição nula
com $2\mathrm p_x$ e com $2\mathrm p_y$ (cada um é antissimétrico em relação a um
plano que contém o eixo, e a combinação é simétrica); eles formam dois pares
não ligantes degenerados. Na flexão no plano $yz$, o $2\mathrm p_y$ passa a
apontar em parte para os hidrogênios: sua sobreposição com $h_1 + h_2$ cresce a
partir de zero, e o nível preenchido que ele carrega é estabilizado (o nível
$3a_1$ da molécula angular). O nível ligante $1b_2$ perde um pouco de
sobreposição e sobe ligeiramente. O saldo é um abaixamento, máximo quando o
nível que desce está preenchido, como é o caso com oito elétrons. Uma molécula
com apenas quatro elétrons de valência, como \ce{BeH2}, tem esse nível vazio e
é linear.
\end{proof}

\section{Metano, amônia e espectros de fotoelétrons}

No metano, os quatro orbitais dos hidrogênios dão a combinação simétrica, que
corresponde ao orbital 2s do carbono, e um conjunto degenerado de três, que
corresponde a seus três orbitais 2p. Quatro \omterm{def:b2:diatomic-mos:bonding}{orbitais ligantes} estão
preenchidos: um baixo ($a_1$) e três degenerados em energia mais alta ($t_2$).
A amônia, com três hidrogênios, conserva como nível ocupado mais alto um
orbital preenchido que aponta para o lado oposto aos hidrogênios, no
nitrogênio: seu par não ligante.

\begin{omfigure}
\begin{tikzpicture}[font=\small]
  \pic (cs) at (0,-1.6) {omlevel={ud}{0.7}}; \node[left] at (-0.45,-1.6) {2s};
  \pic (cp1) at (-0.45,0) {omlevel={u}{0.4}}; \pic (cp2) at (0,0) {omlevel={u}{0.4}}; \pic (cp3) at (0.45,0) {omlevel={empty}{0.4}};
  \node[left] at (-0.7,0) {2p};
  \node at (0,-2.5) {C};
  \pic (ha) at (6,-0.8) {omlevel={ud}{0.7}}; \node[right] at (6.4,-0.8) {$a_1$};
  \pic (ht1) at (5.55,0.4) {omlevel={u}{0.4}}; \pic (ht2) at (6,0.4) {omlevel={u}{0.4}}; \pic (ht3) at (6.45,0.4) {omlevel={empty}{0.4}};
  \node[right] at (6.7,0.4) {$t_2$};
  \node at (6,-1.6) {fragmento \ce{H4}};
  \pic (m1) at (3,-2.4) {omlevel={ud}{0.7}}; \node[omsymlabel, below=0.17cm, inner sep=0.5pt] at (m1-c) {$1a_1$};
  \pic (m2a) at (2.55,-0.9) {omlevel={ud}{0.4}}; \pic (m2b) at (3,-0.9) {omlevel={ud}{0.4}}; \pic (m2c) at (3.45,-0.9) {omlevel={ud}{0.4}};
  \node[omsymlabel, below=0.17cm, inner sep=0.5pt] at (m2b-c) {$1t_2$};
  \pic (m3) at (3,1.8) {omlevel={empty}{0.7}}; \node[omsymlabel, below=0.17cm, inner sep=0.5pt] at (m3-c) {$2a_1$};
  \pic (m4a) at (2.55,2.5) {omlevel={empty}{0.4}}; \pic (m4b) at (3,2.5) {omlevel={empty}{0.4}}; \pic (m4c) at (3.45,2.5) {omlevel={empty}{0.4}};
  \node[omsymlabel, below=0.17cm, inner sep=0.5pt] at (m4b-c) {$2t_2$};
  \draw[omcorr, densely dashed] (cs-r) -- (m1-l) (cp3-r) -- (m2a-l) (cs-r) -- (m3-l) (cp3-r) -- (m4a-l);
  \draw[omcorr, densely dashed] (ha-l) -- (m1-r) (ht1-l) -- (m2c-r) (ha-l) -- (m3-r) (ht1-l) -- (m4c-r);
\end{tikzpicture}

{\small Diagrama de \omterm{def:b2:fragment-orbitals:fragment}{fragmentos} do metano, esquemático. A combinação simétrica
dos quatro hidrogênios corresponde ao orbital 2s do carbono, as três
degeneradas a seus orbitais 2p: dois níveis preenchidos, portanto duas bandas
de fotoelétrons, a mais alta (ionizada primeiro, a \qty{12.61}{eV}) com seis
elétrons.}
% ledger: ie:CH4
\end{omfigure}

\begin{definition}[Espectroscopia de fotoelétrons]\label{def:b2:fragment-orbitals:photoelectron}
A \emph{espectroscopia de fotoelétrons}\index{espectroscopia de fotoeletrons@espectroscopia de fotoelétrons} mede as
energias cinéticas dos elétrons ejetados das moléculas por uma radiação
monocromática de energia conhecida; cada banda do espectro corresponde à
remoção de um elétron de um nível ocupado, e sua posição dá a energia de
ionização desse nível.
\end{definition}

\begin{proposition}[Aproximação de Koopmans]\label{prop:b2:fragment-orbitals:koopmans}
A energia necessária para retirar um elétron de um orbital ocupado é
aproximadamente o oposto da energia desse orbital, $I \approx -\varepsilon$.
\end{proposition}

\admitted

A aproximação despreza a relaxação dos outros elétrons após a ionização; ela
é justificada no volume do 3.º ano. Ela transforma um espectro de fotoelétrons
em um diagrama de orbitais: a água mostra quatro bandas de valência, o metano
duas, de acordo com os diagramas de \omterm{def:b2:fragment-orbitals:fragment}{fragmentos} --- e contra a imagem de dois
pares não ligantes equivalentes e duas ligações equivalentes na água, que
descreve a mesma densidade eletrônica total de outro modo, mas não dá as
energias dos elétrons retirados.

\section{O eteno}

O eteno pode ser construído a partir de dois \omterm{def:b2:fragment-orbitals:fragment}{fragmentos} \ce{CH2}. Cada \ce{CH2}
tem, além dos orbitais de suas duas ligações C--H, um orbital no plano que
aponta para o lado oposto aos hidrogênios (que forma a ligação $\sigma$ C--C
com seu parceiro) e um orbital p perpendicular ao plano. Os dois orbitais p se
combinam em um $\pi$ preenchido e um $\pi^*$ vazio.

\begin{omfigure}
\begin{tikzpicture}[font=\small]
  \begin{scope}
    \pic at (0,0) {omp={0.85}{90}}; \pic at (1.2,0) {omp={0.85}{90}};
    \draw[thick] (0,0) -- (1.2,0);
    \draw[thick] (0,0) -- (-0.5,-0.45) (0,0) -- (-0.55,0.35) (1.2,0) -- (1.7,-0.45) (1.2,0) -- (1.75,0.35);
    \node at (0.6,-1.1) {$\pi$ (preenchido)};
  \end{scope}
  \begin{scope}[xshift=3.6cm]
    \pic at (0,0) {omp={0.85}{90}}; \pic at (1.2,0) {omp={-0.85}{90}};
    \draw[thick] (0,0) -- (1.2,0);
    \draw[thick] (0,0) -- (-0.5,-0.45) (0,0) -- (-0.55,0.35) (1.2,0) -- (1.7,-0.45) (1.2,0) -- (1.75,0.35);
    \node at (0.6,-1.1) {$\pi^*$ (vazio)};
  \end{scope}
  \begin{scope}[xshift=7.4cm]
    \pic at (0,0) {omp={0.85}{90}}; \pic at (1.2,0) {omp={0.85}{0}};
    \draw[thick] (0,0) -- (1.2,0);
    \node at (0.6,-1.1) {torcido de $90^\circ$};
    \node[font=\scriptsize, align=center] at (0.6,1.3) {sobreposição nula};
  \end{scope}
\end{tikzpicture}

{\small O sistema $\pi$ do eteno, esquemático (a molécula vista de perfil, com
os hidrogênios no plano). Os dois orbitais p perpendiculares ao plano se
combinam em fase ($\pi$, preenchido) e fora de fase ($\pi^*$, vazio). Torcer
um \ce{CH2} de $90^\circ$ torna os dois orbitais p perpendiculares: sua
sobreposição, e a ligação $\pi$, se anulam.}
\end{omfigure}

A ligação $\pi$ mantém o eteno plano: torcer uma extremidade afasta os orbitais
p um do outro, sua sobreposição cai com o cosseno do ângulo de torção e, a
$90^\circ$, a ligação $\pi$ desaparece. É preciso fornecer a energia de uma
ligação $\pi$ para torcer a molécula; por isso os isômeros cis e trans dos
alcenos não se interconvertem à temperatura ambiente, enquanto a rotação em
torno de uma ligação simples é livre.

\section{Hiperconjugação}

\begin{definition}[Hiperconjugação]\label{def:b2:fragment-orbitals:hyperconjugation}
A \emph{hiperconjugação}\index{hiperconjugacao@hiperconjugação} é a interação de um orbital
$\sigma$ preenchido de uma ligação C--H ou C--C com um orbital p ou $\pi^*$
vizinho, vazio ou parcialmente preenchido, paralelo a ele.
\end{definition}

\begin{proposition}[Os grupos alquila estabilizam os carbocátions]\label{prop:b2:fragment-orbitals:hyperconjugation}
Um carbocátion é estabilizado por cada ligação C--H ou C--C de um carbono
vizinho que possa se alinhar com seu orbital p vazio: quanto mais ligações
desse tipo, mais estável o cátion; daí a ordem terciário > secundário >
primário > metila.
\end{proposition}

\begin{proof}[Raciocínio]
O orbital p vazio do carbono catiônico e um orbital $\sigma_{\text{C--H}}$
preenchido vizinho formam um sistema de dois níveis com uma grande diferença
$\Delta$ e um pequeno acoplamento $\beta$ (não nulo somente se a ligação for
aproximadamente paralela ao orbital p). Pelo
\cref{thm:b2:diatomic-mos:heteronuclear}, o nível $\sigma$ preenchido é abaixado
de cerca de $\beta^2/\Delta$, e os dois elétrons ganham o dobro disso; o nível
vazio sobe, sem custo. Cada ligação convenientemente alinhada acrescenta sua
própria estabilização; um grupo metila sempre tem uma ligação C--H quase
alinhada, qualquer que seja sua rotação.
\end{proof}

\begin{omfigure}
\begin{tikzpicture}[font=\small]
  \pic at (0,0) {omp={0.9}{90}};
  \node[circle, fill=black, inner sep=1.2pt] at (0,0) {};
  \node[font=\scriptsize, anchor=north west] at (0.08,-0.05) {C$^+$};
  \draw[thick] (0,0) -- (1.4,0);
  \node[circle, fill=black, inner sep=1.2pt] at (1.4,0) {};
  \draw[thick] (1.4,0) -- (1.75,0.85);
  \pic at (1.58,0.42) {omlobe={0.75}{68}};
  \pic at (1.58,0.42) {omlobe={0.45}{248}};
  \node[font=\scriptsize] at (1.95,1.0) {H};
  \draw[thick] (1.4,0) -- (2.0,-0.45) (1.4,0) -- (1.1,-0.6);
  \draw[<->, omProp, thick] (0.35,0.8) .. controls (0.8,1.2) and (1.2,1.2) .. (1.5,0.95);
  \node[font=\scriptsize, align=left, anchor=west] at (2.4,0.4) {lobo $\sigma_{\text{C--H}}$ preenchido paralelo\\ ao orbital p vazio};
  % level scheme
  \begin{scope}[xshift=7.2cm, yshift=-0.4cm]
    \pic (p) at (0,1.4) {omlevel={empty}{0.7}}; \node[left] at (-0.4,1.4) {p vazio};
    \pic (s) at (2.2,-0.8) {omlevel={ud}{0.7}}; \node[right] at (2.6,-0.8) {$\sigma_{\text{C--H}}$};
    \pic (lo) at (1.1,-1.2) {omlevel={ud}{0.7}};
    \pic (hi) at (1.1,1.7) {omlevel={empty}{0.7}};
    \draw[omcorr, densely dashed] (p-r) -- (lo-l) (p-r) -- (hi-l) (s-l) -- (lo-r) (s-l) -- (hi-r);
    \node[font=\scriptsize, anchor=north] at (1.1,-1.35) {estabilizado de $\approx \beta^2/\Delta$};
  \end{scope}
\end{tikzpicture}

{\small \omterm{def:b2:fragment-orbitals:hyperconjugation}{Hiperconjugação} em um carbocátion, esquemática. Um \omterm{def:b2:diatomic-mos:bonding}{orbital ligante} C--H
preenchido no carbono vizinho, paralelo ao orbital p vazio, cede a ele parte
de sua densidade eletrônica; o nível preenchido é estabilizado como em
qualquer interação entre dois níveis de energias diferentes.}
\end{omfigure}

\section{Exercícios}

\begin{exercise}[$\star$]\label{exo:b2:fragment-orbitals:1}
Escreva as duas \omterm{def:b2:fragment-orbitals:symmetry-adapted}{combinações adaptadas à simetria} dos orbitais 1s dos
hidrogênios da água e diga qual delas muda de sinal no plano que contém a
bissetriz do ângulo H--O--H e é perpendicular à molécula.
\end{exercise}

\begin{exercise}[$\star$]\label{exo:b2:fragment-orbitals:2}
Quantos \omterm{def:b2:diatomic-mos:lcao}{orbitais moleculares} de valência tem a água, quantos elétrons de
valência e quantos orbitais de valência preenchidos?
\end{exercise}

\begin{exercise}[$\star$]\label{exo:b2:fragment-orbitals:3}
Por que o espectro de fotoelétrons do metano mostra duas bandas de valência,
uma delas três vezes mais intensa que a outra?
\end{exercise}

\begin{exercise}[$\star$]\label{exo:b2:fragment-orbitals:4}
Por que os seis átomos do eteno estão em um mesmo plano?
\end{exercise}

\begin{exercise}[$\star\star$]\label{exo:b2:fragment-orbitals:5}
A primeira energia de ionização da amônia é \qty{10.07}{eV}, a do metano,
\qty{12.61}{eV}. Que orbital perde o elétron em cada caso, e por que a amônia é
mais fácil de ionizar?
\end{exercise}

\begin{exercise}[$\star\star$]\label{exo:b2:fragment-orbitals:6}
Compare os níveis ocupados da água linear e da água angular e explique por que
uma molécula como \ce{BeH2}, com quatro elétrons de valência, é linear.
\end{exercise}

\begin{exercise}[$\star\star$]\label{exo:b2:fragment-orbitals:7}
Explique, com o resultado dos dois níveis, a ordem de estabilidade dos
carbocátions \ce{CH3+}, \ce{CH3CH2+}, \ce{(CH3)2CH+}, \ce{(CH3)3C+}.
\end{exercise}

\begin{exercise}[$\star\star$]\label{exo:b2:fragment-orbitals:8}
A sobreposição de dois orbitais p torcidos de um ângulo $\theta$ em torno da
ligação é proporcional a $\cos\theta$. Sendo a \omterm{def:b2:diatomic-mos:integrals}{integral de ressonância}
proporcional à sobreposição, como a estabilização $\pi$ depende de $\theta$? Em
que ângulo ela cai à metade?
\end{exercise}

\begin{exercise}[$\star\star$]\label{exo:b2:fragment-orbitals:9}
O borano \ce{BH3} é plano, a amônia, piramidal. Usando os \omterm{def:b2:fragment-orbitals:fragment}{orbitais de fragmento}
de um átomo central e de três hidrogênios, explique a diferença pelos elétrons
de cada um.
\end{exercise}

\begin{exercise}[$\star\star\star$]\label{exo:b2:fragment-orbitals:10}
O íon \ce{H3+} é um triângulo equilátero. Com $\alpha$ e $\beta$ para os orbitais
1s e a sobreposição desprezada, encontre suas três \omterm{def:b2:atomic-orbitals:orbital-energy}{energias orbitais} (a
combinação simétrica tem energia $\alpha + 2\beta$) e explique por que dois
elétrons bastam para ligá-lo.
\end{exercise}

\begin{exercise}[$\star\star\star$]\label{exo:b2:fragment-orbitals:11}
O sistema alila \ce{CH2=CH-CH2} tem três orbitais p paralelos. Adivinhe as
formas de seus três orbitais $\pi$ (sinais em cada carbono) a partir do número
de nós, e qual deles é o nível ocupado mais alto do ânion alila.
\end{exercise}

\begin{exercise}[$\star\star\star$]\label{exo:b2:fragment-orbitals:12}
Construa o sistema $\pi$ do dióxido de carbono a partir dos orbitais 2p
perpendiculares ao eixo O--C--O: quantos orbitais $\pi$, quais são ligantes,
não ligantes, antiligantes, e quantos elétrons eles contêm?
\end{exercise}

\section{Problema: A forma da água}

\begin{problem}[{Problema de fim de semana --- os \omterm{def:b2:fragment-orbitals:fragment}{fragmentos} oxigênio e \ce{H2}, os
orbitais de um H--O--H linear, os níveis que se deslocam quando ele se dobra e
o espectro de fotoelétrons com a aproximação de Koopmans}]
\label{pb:b2:fragment-orbitals:1}
Dados: geometria da água, $r(\ce{O-H}) = \qty{95.8}{pm}$, ângulo
$104.48^\circ$; primeiras energias de ionização: \ce{H2O} \qty{12.62}{eV}, \ce{O}
\qty{13.62}{eV}. A água está no plano $yz$, com $z$ ao longo da bissetriz.
% ledger: geo:H2O.rOH, geo:H2O.aHOH, ie:H2O, ie:O

\medskip
\noindent\textbf{Parte I --- Os \omterm{def:b2:fragment-orbitals:fragment}{fragmentos}.}
\begin{enumerate}
  \item Quantos elétrons de valência tem a água, e de quais átomos?
  \item Escreva as duas \omterm{def:b2:fragment-orbitals:symmetry-adapted}{combinações adaptadas à simetria} dos orbitais 1s dos hidrogênios.
  \item Quais orbitais do oxigênio são simétricos em relação aos dois planos de simetria?
  \item Qual orbital do oxigênio é antissimétrico apenas em relação ao plano $xz$?
  \item Qual orbital do oxigênio não tem parceiro entre as combinações dos hidrogênios?
  \item Quantos \omterm{def:b2:diatomic-mos:lcao}{orbitais moleculares} resultam, e quantos estão preenchidos?
\end{enumerate}

\medskip
\noindent\textbf{Parte II --- Um H--O--H linear.}
\begin{enumerate}[resume]
  \item Em uma molécula linear ao longo de $z$, qual combinação interage com o
        $2\mathrm p_z$ do oxigênio?
  \item Qual interage com o 2s do oxigênio?
  \item Quais orbitais do oxigênio permanecem não ligantes, e com que degenerescência?
  \item Preencha os níveis com os oito elétrons de valência.
  \item A água linear seria \omterm{def:b2:diatomic-mos:magnetism}{paramagnética}?
  \item Onde ficaria seu nível ocupado mais alto, comparado com um orbital 2p
        do oxigênio?
\end{enumerate}

\medskip
\noindent\textbf{Parte III --- A flexão.}
\begin{enumerate}[resume]
  \item A molécula linear, ao longo de $z$, é dobrada no plano $yz$. Qual
        \omterm{def:b2:diatomic-mos:bonding}{orbital não ligante} começa a se sobrepor a $h_1 + h_2$?
  \item O que acontece com o nível preenchido que ele carrega?
  \item Qual nível ligante sobe ligeiramente, e por quê?
  \item Conclua: por que a água é angular?
  \item Compare o ângulo medido com $90^\circ$ (ligação por orbitais p puros) e com $109.5^\circ$.
  \item Qual \omterm{def:b2:diatomic-mos:bonding}{orbital não ligante} não é afetado pela flexão?
\end{enumerate}

\medskip
\noindent\textbf{Parte IV --- O espectro de fotoelétrons.}
\begin{enumerate}[resume]
  \item Quantas bandas de valência o espectro mostra?
  \item Pela aproximação de Koopmans, de qual orbital vem a primeira banda?
  \item Por que essa banda é estreita, enquanto as bandas dos \omterm{def:b2:diatomic-mos:bonding}{orbitais ligantes} são largas?
  \item Compare a primeira energia de ionização da água com a do átomo de
        oxigênio e explique a diferença.
  \item O que acontece com os ``dois pares não ligantes equivalentes'' da estrutura de Lewis?
  \item Dê o número de bandas de ionização de valência da água e a energia da
        mais baixa.
\end{enumerate}
\end{problem}
